\item \textbf{Clasificación de número entero}.

\begin{tcolorbox}[colback=blue!2!white,colframe=blue!55!black,title={Enunciado}]
\begin{itemize}
	\item Escriba un programa en \texttt{Python} que solicite al usuario ingresar un número entero.
	\item El programa debe analizar el valor ingresado y determinar si corresponde a un número par o impar. Para lograr el objetivo, divida el número ingresado por 2, y evalúe el resto utilizando la operación módulo \texttt{ \% }
\item Para realizar la clasificación, utilice únicamente estructuras condicionales (\texttt{if}, \texttt{elif}, \texttt{else}).
\item Finalmente, el programa debe mostrar un mensaje indicando si el número ingresado es un par o impar.
\end{itemize}
\end{tcolorbox}

\begin{solution}

\subsubsection*{Idea}

El programa debe solicitar al usuario un número entero.  
Luego debe verificar si el número es divisible por 2.

Para realizar esta verificación se utiliza la operación módulo \texttt{\%}.  
El operador módulo devuelve el resto de una división.

Por lo tanto:

\begin{itemize}
\item Si el resto de dividir el número por 2 es igual a 0, el número es \textbf{par}.
\item Si el resto es diferente de 0, el número es \textbf{impar}.
\end{itemize}

Esta clasificación se puede realizar utilizando estructuras condicionales.

\subsubsection*{Implementación}

\begin{pythonjupyter}
numero = int(input("Ingrese un número entero: "))

if numero % 2 == 0:
    print("El número ingresado es par.")
else:
    print("El número ingresado es impar.")
\end{pythonjupyter}

\subsubsection*{Explicación del algoritmo}

Primero se solicita al usuario un número entero:

\begin{pythonjupyter}
numero = int(input("Ingrese un número entero: "))
\end{pythonjupyter}

La función \texttt{int()} convierte el valor ingresado a un número entero.

Luego se evalúa la condición utilizando el operador módulo:

\begin{pythonjupyter}
if numero % 2 == 0:
\end{pythonjupyter}

Aquí ocurre lo siguiente:

\begin{itemize}
\item \texttt{numero \% 2} calcula el resto de dividir el número por 2.
\item Si el resto es 0, el número es divisible por 2.
\item Todo número divisible por 2 es un número par.
\end{itemize}

Ejemplo:

\begin{itemize}
\item Si el usuario ingresa \texttt{8}  
      entonces \texttt{8 \% 2 = 0}.  
      Por lo tanto el número es par.

\item Si el usuario ingresa \texttt{7}  
      entonces \texttt{7 \% 2 = 1}.  
      Por lo tanto el número es impar.
\end{itemize}

Dependiendo del resultado de la condición, el programa imprime el mensaje correspondiente mediante \texttt{print()}.

\subsubsection*{Resultado}

El programa permite clasificar correctamente si el número ingresado es par o impar utilizando la operación módulo y estructuras condicionales.

\end{solution}
